\chapter{Unit 4: Fourier Series and Harmonic Expansions}

\section{Section A: Multiple Choice Practice Questions (60 Curated Problems)}

\noindent\textbf{Q1. What is the primary purpose of a Fourier series?} \hfill [Tag: introduction and Euler's formulae | Easy]
\begin{enumerate}[label=(\alph*),topsep=2pt,itemsep=1pt,partopsep=0pt,parsep=0pt]
    \item {{To calculate the limits of continuous functions}}
    \item {{To find the roots of a polynomial equation}}
    \item {{To solve differential equations directly}}
    \item {{To represent a periodic function as an infinite sum of sines and cosines}}
\end{enumerate}
\vspace{2mm}

\noindent\textbf{Q2. In the standard Fourier series of f ( x ) over [ - $\pi$ , $\pi$ ] , what does a 0 represent?} \hfill [Tag: introduction and Euler's formulae | Easy]
\begin{enumerate}[label=(\alph*),topsep=2pt,itemsep=1pt,partopsep=0pt,parsep=0pt]
    \item {{The amplitude of the sine waves}}
    \item {{The fundamental frequency}}
    \item {{Twice the average value of the function over one period}}
    \item {{The maximum value of the function}}
\end{enumerate}
\vspace{2mm}

\noindent\textbf{Q3. Which of the following represents Euler's formula for the Fourier coefficient a n over the interval [ - $\pi$ , $\pi$ ] ?} \hfill [Tag: introduction and Euler's formulae | Easy]
\begin{enumerate}[label=(\alph*),topsep=2pt,itemsep=1pt,partopsep=0pt,parsep=0pt]
    \item {{$\pi$}}
    \item {{$\pi$}}
    \item {{2 $\pi$}}
    \item {{$\int$ - $\pi$}}
\end{enumerate}
\vspace{2mm}

\noindent\textbf{Q4. What is the formula for the Fourier coefficient b n for a function f ( x ) on [ - $\pi$ , $\pi$ ] ?} \hfill [Tag: introduction and Euler's formulae | Easy]
\begin{enumerate}[label=(\alph*),topsep=2pt,itemsep=1pt,partopsep=0pt,parsep=0pt]
    \item {{$\int$ - $\pi$}}
    \item {{2 $\pi$}}
    \item {{$\pi$}}
    \item {{$\pi$}}
\end{enumerate}
\vspace{2mm}

\noindent\textbf{Q5. Which set of conditions determines whether a function can be represented by a Fourier series?} \hfill [Tag: conditions for a Fourier expansion and functions having points of discontinuity | Easy]
\begin{enumerate}[label=(\alph*),topsep=2pt,itemsep=1pt,partopsep=0pt,parsep=0pt]
    \item {{Laplace's criteria}}
    \item {{Cauchy-Riemann equations}}
    \item {{Euler's conditions}}
    \item {{Dirichlet's conditions}}
\end{enumerate}
\vspace{2mm}

\noindent\textbf{Q6. According to Dirichlet's conditions, the number of finite discontinuities in one period must be:} \hfill [Tag: conditions for a Fourier expansion and functions having points of discontinuity | Easy]
\begin{enumerate}[label=(\alph*),topsep=2pt,itemsep=1pt,partopsep=0pt,parsep=0pt]
    \item {{Finite}}
    \item {{Zero}}
    \item {{Exactly one}}
    \item {{Infinite}}
\end{enumerate}
\vspace{2mm}

\noindent\textbf{Q7. If a function f ( x ) has a point of finite discontinuity at x = c , to what value does the Fourier series converge at x = c ?} \hfill [Tag: conditions for a Fourier expansion and functions having points of discontinuity | Easy]
\begin{enumerate}[label=(\alph*),topsep=2pt,itemsep=1pt,partopsep=0pt,parsep=0pt]
    \item {{Zero}}
    \item {{The average of the left-hand and right-hand limits: 2}}
    \item {{The left-hand limit: f ( c -}}
    \item {{The right-hand limit: f ( c +}}
\end{enumerate}
\vspace{2mm}

\noindent\textbf{Q8. Which of the following violates Dirichlet's conditions for a Fourier series?} \hfill [Tag: conditions for a Fourier expansion and functions having points of discontinuity | Easy]
\begin{enumerate}[label=(\alph*),topsep=2pt,itemsep=1pt,partopsep=0pt,parsep=0pt]
    \item {{Having an infinite number of maxima and minima in a single period}}
    \item {{Having a finite number of jump discontinuities}}
    \item {{Being absolutely integrable over a period}}
    \item {{Being a single-valued function}}
\end{enumerate}
\vspace{2mm}

\noindent\textbf{Q9. When expanding a function f ( x ) defined over the interval [ - l , l ] , what is the argument of the sine and cosine terms in the Fourier series?} \hfill [Tag: change of interval | Easy]
\begin{enumerate}[label=(\alph*),topsep=2pt,itemsep=1pt,partopsep=0pt,parsep=0pt]
    \item {{n$\pi$ x}}
    \item {{l}}
    \item {{l}}
    \item {{n x}}
\end{enumerate}
\vspace{2mm}

\noindent\textbf{Q10. What is the period of the terms cos ( l n$\pi$ x ) and sin ( l n$\pi$ x ) ?} \hfill [Tag: change of interval | Easy]
\begin{enumerate}[label=(\alph*),topsep=2pt,itemsep=1pt,partopsep=0pt,parsep=0pt]
    \item {{l}}
    \item {{n}}
    \item {{2 $\pi$}}
    \item {{2 l}}
\end{enumerate}
\vspace{2mm}

\noindent\textbf{Q11. For a function defined over [ 0 , 2 L ] , what is the multiplier in front of the integral for the Fourier coefficient a 0 ?} \hfill [Tag: change of interval | Easy]
\begin{enumerate}[label=(\alph*),topsep=2pt,itemsep=1pt,partopsep=0pt,parsep=0pt]
    \item {{L}}
    \item {{$\pi$}}
    \item {{2 L}}
    \item {{L}}
\end{enumerate}
\vspace{2mm}

\noindent\textbf{Q12. If a periodic function f ( x ) is an even function, which of its Fourier coefficients will evaluate to zero?} \hfill [Tag: even and odd functions | Easy]
\begin{enumerate}[label=(\alph*),topsep=2pt,itemsep=1pt,partopsep=0pt,parsep=0pt]
    \item {{All a n}}
    \item {{a 0}}
    \item {{Both a n}}
    \item {{All b n}}
\end{enumerate}
\vspace{2mm}

\noindent\textbf{Q13. If f ( - x ) = - f ( x ) for all x , what kind of function is f ( x ) ?} \hfill [Tag: even and odd functions | Easy]
\begin{enumerate}[label=(\alph*),topsep=2pt,itemsep=1pt,partopsep=0pt,parsep=0pt]
    \item {{Odd function}}
    \item {{Even function}}
    \item {{Neither even nor odd}}
    \item {{Constant function}}
\end{enumerate}
\vspace{2mm}

\noindent\textbf{Q14. The Fourier series of an odd function f ( x ) over the interval [ - $\pi$ , $\pi$ ] consists entirely of:} \hfill [Tag: even and odd functions | Easy]
\begin{enumerate}[label=(\alph*),topsep=2pt,itemsep=1pt,partopsep=0pt,parsep=0pt]
    \item {{Both sine and cosine terms}}
    \item {{Sine terms}}
    \item {{Cosine terms}}
    \item {{Constant terms}}
\end{enumerate}
\vspace{2mm}

\noindent\textbf{Q15. What is the value of $\int$ - L L f ( x ) d x if f ( x ) is an odd function?} \hfill [Tag: even and odd functions | Easy]
\begin{enumerate}[label=(\alph*),topsep=2pt,itemsep=1pt,partopsep=0pt,parsep=0pt]
    \item {{0}}
    \item {{2 L}}
    \item {{2 $\int$ 0}}
    \item {{L}}
\end{enumerate}
\vspace{2mm}

\noindent\textbf{Q16. Which of the following functions is an even function?} \hfill [Tag: even and odd functions | Easy]
\begin{enumerate}[label=(\alph*),topsep=2pt,itemsep=1pt,partopsep=0pt,parsep=0pt]
    \item {{f ( x ) =}}
    \item {{f ( x ) =}}
    \item {{f ( x ) =}}
    \item {{f ( x ) =}}
\end{enumerate}
\vspace{2mm}

\noindent\textbf{Q17. A half range cosine series for a function f ( x ) defined on [ 0 , L ] is obtained by extending f ( x ) to [ - L , L ] as an:} \hfill [Tag: half range series | Easy]
\begin{enumerate}[label=(\alph*),topsep=2pt,itemsep=1pt,partopsep=0pt,parsep=0pt]
    \item {{Periodic function of period L}}
    \item {{Odd function}}
    \item {{Identity function}}
    \item {{Even function}}
\end{enumerate}
\vspace{2mm}

\noindent\textbf{Q18. To find the half range sine series of f ( x ) over [ 0 , $\pi$ ] , we assume the extended function F ( x ) on [ - $\pi$ , $\pi$ ] satisfies:} \hfill [Tag: half range series | Easy]
\begin{enumerate}[label=(\alph*),topsep=2pt,itemsep=1pt,partopsep=0pt,parsep=0pt]
    \item {{F ( - x ) =}}
    \item {{F ( - x ) =}}
    \item {{F ( x +}}
    \item {{F ( x ) =}}
\end{enumerate}
\vspace{2mm}

\noindent\textbf{Q19. For a half range sine series of f ( x ) over [ 0 , L ] , what is the formula for the coefficient b n ?} \hfill [Tag: half range series | Easy]
\begin{enumerate}[label=(\alph*),topsep=2pt,itemsep=1pt,partopsep=0pt,parsep=0pt]
    \item {{L}}
    \item {{L}}
    \item {{L}}
    \item {{$\int$ 0}}
\end{enumerate}
\vspace{2mm}

\noindent\textbf{Q20. Which coefficients are always zero in a half range sine series?} \hfill [Tag: half range series | Easy]
\begin{enumerate}[label=(\alph*),topsep=2pt,itemsep=1pt,partopsep=0pt,parsep=0pt]
    \item {{a 0}}
    \item {{a n}}
    \item {{a 0}}
    \item {{b n}}
\end{enumerate}
\vspace{2mm}

\noindent\textbf{Q21. What is the value of the Fourier coefficient a 0 for the function f ( x ) = x in the interval ( - $\pi$ , $\pi$ ) ?} \hfill [Tag: introduction and Euler's formulae | Medium]
\begin{enumerate}[label=(\alph*),topsep=2pt,itemsep=1pt,partopsep=0pt,parsep=0pt]
    \item {{$\pi$}}
    \item {{0}}
    \item {{- $\pi$}}
    \item {{1}}
\end{enumerate}
\vspace{2mm}

\noindent\textbf{Q22. According to Dirichlet's conditions, a function f ( x ) can be expanded into a Fourier series if it is:} \hfill [Tag: conditions for a Fourier expansion and functions having points of discontinuity | Medium]
\begin{enumerate}[label=(\alph*),topsep=2pt,itemsep=1pt,partopsep=0pt,parsep=0pt]
    \item {{Continuous everywhere}}
    \item {{Infinite at certain points in the interval}}
    \item {{Single-valued, finite, and has a finite number of finite discontinuities}}
    \item {{Differentiable everywhere}}
\end{enumerate}
\vspace{2mm}

\noindent\textbf{Q23. At a point of finite discontinuity x = c , the Fourier series of f ( x ) converges to:} \hfill [Tag: conditions for a Fourier expansion and functions having points of discontinuity | Medium]
\begin{enumerate}[label=(\alph*),topsep=2pt,itemsep=1pt,partopsep=0pt,parsep=0pt]
    \item {{Infinity}}
    \item {{f ( c )}}
    \item {{2}}
    \item {{0}}
\end{enumerate}
\vspace{2mm}

\noindent\textbf{Q24. If f ( x ) = | x | in the interval ( - $\pi$ , $\pi$ ) , which Fourier coefficients are zero?} \hfill [Tag: even and odd functions | Medium]
\begin{enumerate}[label=(\alph*),topsep=2pt,itemsep=1pt,partopsep=0pt,parsep=0pt]
    \item {{a n}}
    \item {{a 0}}
    \item {{Both a n}}
    \item {{b n}}
\end{enumerate}
\vspace{2mm}

\noindent\textbf{Q25. The Euler formula for the Fourier coefficient a n in the interval ( 0 , 2 $\pi$ ) is given by:} \hfill [Tag: introduction and Euler's formulae | Medium]
\begin{enumerate}[label=(\alph*),topsep=2pt,itemsep=1pt,partopsep=0pt,parsep=0pt]
    \item {{$\pi$}}
    \item {{$\pi$}}
    \item {{$\pi$}}
    \item {{2 $\pi$}}
\end{enumerate}
\vspace{2mm}

\noindent\textbf{Q26. For a function f ( x ) defined in the interval ( - l , l ) , the argument of the trigonometric functions in the Fourier series is:} \hfill [Tag: change of interval | Medium]
\begin{enumerate}[label=(\alph*),topsep=2pt,itemsep=1pt,partopsep=0pt,parsep=0pt]
    \item {{n x}}
    \item {{l}}
    \item {{2 l}}
    \item {{l}}
\end{enumerate}
\vspace{2mm}

\noindent\textbf{Q27. If a function f ( x ) is defined over ( 0 , 2 l ) , what is the formula for the Fourier coefficient b n ?} \hfill [Tag: change of interval | Medium]
\begin{enumerate}[label=(\alph*),topsep=2pt,itemsep=1pt,partopsep=0pt,parsep=0pt]
    \item {{l}}
    \item {{l}}
    \item {{l}}
    \item {{2 l}}
\end{enumerate}
\vspace{2mm}

\noindent\textbf{Q28. Which of the following functions will have a Fourier series consisting only of sine terms in ( - $\pi$ , $\pi$ ) ?} \hfill [Tag: even and odd functions | Medium]
\begin{enumerate}[label=(\alph*),topsep=2pt,itemsep=1pt,partopsep=0pt,parsep=0pt]
    \item {{f ( x ) =}}
    \item {{f ( x ) =}}
    \item {{f ( x ) =}}
    \item {{f ( x ) =}}
\end{enumerate}
\vspace{2mm}

\noindent\textbf{Q29. To express f ( x ) = x as a half-range cosine series in the interval ( 0 , $\pi$ ) , we extend f ( x ) to ( - $\pi$ , $\pi$ ) such that the extended function is:} \hfill [Tag: half range series | Medium]
\begin{enumerate}[label=(\alph*),topsep=2pt,itemsep=1pt,partopsep=0pt,parsep=0pt]
    \item {{Zero in ( - $\pi$ , 0 )}}
    \item {{Odd}}
    \item {{Even}}
    \item {{Periodic with period $\pi$}}
\end{enumerate}
\vspace{2mm}

\noindent\textbf{Q30. What is the formula for b n in the half-range sine series of f ( x ) over ( 0 , l ) ?} \hfill [Tag: half range series | Medium]
\begin{enumerate}[label=(\alph*),topsep=2pt,itemsep=1pt,partopsep=0pt,parsep=0pt]
    \item {{l}}
    \item {{0}}
    \item {{l}}
    \item {{l}}
\end{enumerate}
\vspace{2mm}

\noindent\textbf{Q31. If f ( x ) = sin x in ( - $\pi$ , $\pi$ ) , what are the values of a 1 and b 1 respectively?} \hfill [Tag: introduction and Euler's formulae | Medium]
\begin{enumerate}[label=(\alph*),topsep=2pt,itemsep=1pt,partopsep=0pt,parsep=0pt]
    \item {{1 , 1}}
    \item {{1 , 0}}
    \item {{0 , 1}}
    \item {{0 , 0}}
\end{enumerate}
\vspace{2mm}

\noindent\textbf{Q32. Consider f ( x ) = 0 for - $\pi$ < x < 0 and f ( x ) = $\pi$ for 0 < x < $\pi$ . At x = 0 , what does the Fourier series sum to?} \hfill [Tag: conditions for a Fourier expansion and functions having points of discontinuity | Medium]
\begin{enumerate}[label=(\alph*),topsep=2pt,itemsep=1pt,partopsep=0pt,parsep=0pt]
    \item {{0}}
    \item {{Undefined}}
    \item {{2}}
    \item {{$\pi$}}
\end{enumerate}
\vspace{2mm}

\noindent\textbf{Q33. If a function is defined as f ( x ) = x 2 for - $\pi$ < x < $\pi$ , which statement is true about its Fourier coefficients?} \hfill [Tag: even and odd functions | Medium]
\begin{enumerate}[label=(\alph*),topsep=2pt,itemsep=1pt,partopsep=0pt,parsep=0pt]
    \item {{a n}}
    \item {{a 0}}
    \item {{All coefficients are non-zero}}
    \item {{b n}}
\end{enumerate}
\vspace{2mm}

\noindent\textbf{Q34. In the half-range cosine series of a function f ( x ) over the interval ( 0 , $\pi$ ) , the value of b n is always:} \hfill [Tag: half range series | Medium]
\begin{enumerate}[label=(\alph*),topsep=2pt,itemsep=1pt,partopsep=0pt,parsep=0pt]
    \item {{0}}
    \item {{Dependent on f ( x )}}
    \item {{$\pi$}}
    \item {{$\pi$}}
\end{enumerate}
\vspace{2mm}

\noindent\textbf{Q35. If a function has a period of T = 2 l , the fundamental frequency of the corresponding Fourier series is:} \hfill [Tag: change of interval | Medium]
\begin{enumerate}[label=(\alph*),topsep=2pt,itemsep=1pt,partopsep=0pt,parsep=0pt]
    \item {{l}}
    \item {{l}}
    \item {{2 l}}
    \item {{l}}
\end{enumerate}
\vspace{2mm}

\noindent\textbf{Q36. Which integral represents the orthogonality property used to derive Euler's formula for a n ( n  = m ) in the interval ( - $\pi$ , $\pi$ ) ?} \hfill [Tag: introduction and Euler's formulae | Medium]
\begin{enumerate}[label=(\alph*),topsep=2pt,itemsep=1pt,partopsep=0pt,parsep=0pt]
    \item {{$\int$ - $\pi$}}
    \item {{All of the above}}
    \item {{$\int$ - $\pi$}}
    \item {{$\int$ - $\pi$}}
\end{enumerate}
\vspace{2mm}

\noindent\textbf{Q37. Which of the following functions does NOT satisfy Dirichlet's conditions in the interval ( - $\pi$ , $\pi$ ) ?} \hfill [Tag: conditions for a Fourier expansion and functions having points of discontinuity | Medium]
\begin{enumerate}[label=(\alph*),topsep=2pt,itemsep=1pt,partopsep=0pt,parsep=0pt]
    \item {{f ( x ) =}}
    \item {{f ( x ) =}}
    \item {{f ( x ) =}}
    \item {{f ( x ) =}}
\end{enumerate}
\vspace{2mm}

\noindent\textbf{Q38. The product of an even function and an odd function is:} \hfill [Tag: even and odd functions | Medium]
\begin{enumerate}[label=(\alph*),topsep=2pt,itemsep=1pt,partopsep=0pt,parsep=0pt]
    \item {{Neither even nor odd}}
    \item {{Odd}}
    \item {{Zero}}
    \item {{Even}}
\end{enumerate}
\vspace{2mm}

\noindent\textbf{Q39. To expand f ( x ) = x 2 as a half-range sine series in ( 0 , 2 ) , what extension is applied?} \hfill [Tag: half range series | Medium]
\begin{enumerate}[label=(\alph*),topsep=2pt,itemsep=1pt,partopsep=0pt,parsep=0pt]
    \item {{f ( x ) =}}
    \item {{f ( x ) =}}
    \item {{f ( x ) =}}
    \item {{f ( x ) =}}
\end{enumerate}
\vspace{2mm}

\noindent\textbf{Q40. If f ( x ) is periodic with period 4 and f ( x ) = 1 - x for 0 < x < 4 , the Fourier coefficient a 0 is evaluated as:} \hfill [Tag: change of interval | Medium]
\begin{enumerate}[label=(\alph*),topsep=2pt,itemsep=1pt,partopsep=0pt,parsep=0pt]
    \item {{2}}
    \item {{4}}
    \item {{4}}
    \item {{$\int$ 0}}
\end{enumerate}
\vspace{2mm}

\noindent\textbf{Q41. Let f ( x ) be an absolutely integrable function on [ - $\pi$ , $\pi$ ] . According to the Riemann-Lebesgue Lemma, what is the behavior of the Fourier coefficients a n and b n as n $\to$ $\infty$ ?} \hfill [Tag: introduction and Euler's formulae | Hard]
\begin{enumerate}[label=(\alph*),topsep=2pt,itemsep=1pt,partopsep=0pt,parsep=0pt]
    \item {{They grow proportionally to n}}
    \item {{They oscillate infinitely without converging.}}
    \item {{They approach 0}}
    \item {{They approach 1}}
\end{enumerate}
\vspace{2mm}

\noindent\textbf{Q42. If the Fourier series of a continuously differentiable periodic function f ( x ) is 2 a 0 + $\sum$ n = 1 $\infty$ ( a n cos n x + b n sin n x ) , what are the Fourier coefficients A n and B n for its derivative f ' ( x ) ?} \hfill [Tag: introduction and Euler's formulae | Hard]
\begin{enumerate}[label=(\alph*),topsep=2pt,itemsep=1pt,partopsep=0pt,parsep=0pt]
    \item {{A n}}
    \item {{A n}}
    \item {{A n}}
    \item {{A n}}
\end{enumerate}
\vspace{2mm}

\noindent\textbf{Q43. Let f ( x ) = x 2 on [ - $\pi$ , $\pi$ ] extended periodically. Which of the following best represents the rate of decay of its Fourier coefficients a n ?} \hfill [Tag: introduction and Euler's formulae | Hard]
\begin{enumerate}[label=(\alph*),topsep=2pt,itemsep=1pt,partopsep=0pt,parsep=0pt]
    \item {{Exponential decay}}
    \item {{O ( 1/ n 3}}
    \item {{O ( 1/ n 2}}
    \item {{O ( 1/ n )}}
\end{enumerate}
\vspace{2mm}

\noindent\textbf{Q44. Consider a periodic function f ( x ) of period 2 $\pi$ that satisfies Parseval's identity. If a 0 = 4 and a n = 2 n 1 , b n = 0 for n $\ge$ 1 , what is the value of $\pi$ 1 $\int$ - $\pi$ $\pi$ [ f ( x ) ] 2 d x ?} \hfill [Tag: introduction and Euler's formulae | Hard]
\begin{enumerate}[label=(\alph*),topsep=2pt,itemsep=1pt,partopsep=0pt,parsep=0pt]
    \item {{3}}
    \item {{3}}
    \item {{3}}
    \item {{9}}
\end{enumerate}
\vspace{2mm}

\noindent\textbf{Q45. Which of the following functions violates Dirichlet's conditions for Fourier series expansion on [ - $\pi$ , $\pi$ ] ?} \hfill [Tag: conditions for a Fourier expansion and functions having points of discontinuity | Hard]
\begin{enumerate}[label=(\alph*),topsep=2pt,itemsep=1pt,partopsep=0pt,parsep=0pt]
    \item {{f ( x ) =}}
    \item {{f ( x ) =}}
    \item {{f ( x ) =}}
    \item {{f ( x ) =}}
\end{enumerate}
\vspace{2mm}

\noindent\textbf{Q46. Let f ( x ) be a periodic function with period 2 $\pi$ . If f ( x ) = 0 for - $\pi$ < x < 0 and f ( x ) = 1 for 0 < x < $\pi$ , to what value does the Fourier series of f ( x ) converge at x = 2 $\pi$ ?} \hfill [Tag: conditions for a Fourier expansion and functions having points of discontinuity | Hard]
\begin{enumerate}[label=(\alph*),topsep=2pt,itemsep=1pt,partopsep=0pt,parsep=0pt]
    \item {{1/2}}
    \item {{1}}
    \item {{0}}
    \item {{The series does not converge}}
\end{enumerate}
\vspace{2mm}

\noindent\textbf{Q47. What is the phenomenon called where the partial sums of a Fourier series exhibit an overshoot of approximately 9\% near a jump discontinuity, regardless of the number of terms?} \hfill [Tag: conditions for a Fourier expansion and functions having points of discontinuity | Hard]
\begin{enumerate}[label=(\alph*),topsep=2pt,itemsep=1pt,partopsep=0pt,parsep=0pt]
    \item {{Parseval's Error}}
    \item {{Runge's Phenomenon}}
    \item {{Dirichlet Overshoot}}
    \item {{Gibbs Phenomenon}}
\end{enumerate}
\vspace{2mm}

\noindent\textbf{Q48. A function f ( x ) is defined on [ - $\pi$ , $\pi$ ] and satisfies Dirichlet's conditions. If f ( x ) is continuous everywhere except at x = c where it has a jump discontinuity, the Fourier series at x = c converges to 2 f ( c + ) + f ( c - ) . Which theorem formally guarantees this pointwise convergence?} \hfill [Tag: conditions for a Fourier expansion and functions having points of discontinuity | Hard]
\begin{enumerate}[label=(\alph*),topsep=2pt,itemsep=1pt,partopsep=0pt,parsep=0pt]
    \item {{Fejér's Theorem}}
    \item {{Weierstrass Approximation Theorem}}
    \item {{Dirichlet's Theorem}}
    \item {{Plancherel Theorem}}
\end{enumerate}
\vspace{2mm}

\noindent\textbf{Q49. For a function f ( x ) defined on [ - L , L ] , the Fourier series is given in terms of cos ( L n$\pi$ x ) and sin ( L n$\pi$ x ) . What is the constant coefficient a 0 in this expansion?} \hfill [Tag: change of interval | Hard]
\begin{enumerate}[label=(\alph*),topsep=2pt,itemsep=1pt,partopsep=0pt,parsep=0pt]
    \item {{L}}
    \item {{2 L}}
    \item {{L}}
    \item {{$\int$ - L}}
\end{enumerate}
\vspace{2mm}

\noindent\textbf{Q50. If f ( x ) = x 2 on the interval [ 0 , 2 L ] , what is the appropriate argument for the trigonometric basis functions to form an orthogonal set over this interval?} \hfill [Tag: change of interval | Hard]
\begin{enumerate}[label=(\alph*),topsep=2pt,itemsep=1pt,partopsep=0pt,parsep=0pt]
    \item {{n$\pi$ x}}
    \item {{L}}
    \item {{L}}
    \item {{2 L}}
\end{enumerate}
\vspace{2mm}

\noindent\textbf{Q51. Let f ( x ) be defined on [ 0 , c ] . To find a Fourier series of f ( x ) using standard basis functions without extending the function's domain symmetrically, what must be the fundamental period T of the expansion?} \hfill [Tag: change of interval | Hard]
\begin{enumerate}[label=(\alph*),topsep=2pt,itemsep=1pt,partopsep=0pt,parsep=0pt]
    \item {{c}}
    \item {{2 c}}
    \item {{$\pi$ / c}}
    \item {{c /2}}
\end{enumerate}
\vspace{2mm}

\noindent\textbf{Q52. Consider f ( x ) = e - x on [ - l , l ] . The Fourier coefficient a n will contain a factor of l 2 + ( n$\pi$ ) 2 1 . What is the integral $\int$ - l l e - x cos ( l n$\pi$ x ) d x directly proportional to?} \hfill [Tag: change of interval | Hard]
\begin{enumerate}[label=(\alph*),topsep=2pt,itemsep=1pt,partopsep=0pt,parsep=0pt]
    \item {{( - 1 ) n}}
    \item {{sinh ( l )}}
    \item {{cosh ( l )}}
    \item {{( - 1 ) n}}
\end{enumerate}
\vspace{2mm}

\noindent\textbf{Q53. A periodic function f ( t ) with period T satisfies f ( t ) = - f ( t + T /2 ) for all t . What implication does this 'half-wave symmetry' have on its Fourier series?} \hfill [Tag: even and odd functions | Hard]
\begin{enumerate}[label=(\alph*),topsep=2pt,itemsep=1pt,partopsep=0pt,parsep=0pt]
    \item {{The series contains only even harmonics.}}
    \item {{The series contains only odd harmonics.}}
    \item {{The series contains only cosine terms.}}
    \item {{The series contains only sine terms.}}
\end{enumerate}
\vspace{2mm}

\noindent\textbf{Q54. Let f ( x ) = x sin ( x ) for - $\pi$ < x < $\pi$ . Which statement perfectly describes its Fourier coefficients?} \hfill [Tag: even and odd functions | Hard]
\begin{enumerate}[label=(\alph*),topsep=2pt,itemsep=1pt,partopsep=0pt,parsep=0pt]
    \item {{Both a n}}
    \item {{b n}}
    \item {{a n}}
    \item {{a n}}
\end{enumerate}
\vspace{2mm}

\noindent\textbf{Q55. Consider a function f ( x ) defined on [ - $\pi$ , $\pi$ ] . If f ( x ) is both even and possesses half-wave symmetry, which of the following is true regarding its Fourier coefficients?} \hfill [Tag: even and odd functions | Hard]
\begin{enumerate}[label=(\alph*),topsep=2pt,itemsep=1pt,partopsep=0pt,parsep=0pt]
    \item {{a n}}
    \item {{a n}}
    \item {{a n}}
    \item {{a n}}
\end{enumerate}
\vspace{2mm}

\noindent\textbf{Q56. What is the result of decomposing the function f ( x ) = e i x on [ - $\pi$ , $\pi$ ] into its even and odd parts?} \hfill [Tag: even and odd functions | Hard]
\begin{enumerate}[label=(\alph*),topsep=2pt,itemsep=1pt,partopsep=0pt,parsep=0pt]
    \item {{Even part is sin x}}
    \item {{Even part is cos x}}
    \item {{Even part is cosh x}}
    \item {{Even part is 1}}
\end{enumerate}
\vspace{2mm}

\noindent\textbf{Q57. To find the half-range sine series of f ( x ) = x on ( 0 , L ) , we analytically extend f ( x ) to ( - L , 0 ) . What is the extended function F ( x ) on ( - L , 0 ) ?} \hfill [Tag: half range series | Hard]
\begin{enumerate}[label=(\alph*),topsep=2pt,itemsep=1pt,partopsep=0pt,parsep=0pt]
    \item {{F ( x ) =}}
    \item {{F ( x ) =}}
    \item {{F ( x ) =}}
    \item {{F ( x ) =}}
\end{enumerate}
\vspace{2mm}

\noindent\textbf{Q58. Find the half-range cosine series for f ( x ) = x in ( 0 , $\pi$ ) . What is the general expression for a n (for n $\ge$ 1 )?} \hfill [Tag: half range series | Hard]
\begin{enumerate}[label=(\alph*),topsep=2pt,itemsep=1pt,partopsep=0pt,parsep=0pt]
    \item {{$\pi$ n 2}}
    \item {{0}}
    \item {{$\pi$n}}
    \item {{$\pi$ n 2}}
\end{enumerate}
\vspace{2mm}

\noindent\textbf{Q59. Using the half-range cosine series of f ( x ) = x on ( 0 , $\pi$ ) , which infinite series sum can be deduced by evaluating the series at x = 0 ?} \hfill [Tag: half range series | Hard]
\begin{enumerate}[label=(\alph*),topsep=2pt,itemsep=1pt,partopsep=0pt,parsep=0pt]
    \item {{$\sum$ k = 1}}
    \item {{$\sum$ k = 1}}
    \item {{$\sum$ k = 1}}
    \item {{$\sum$ k = 1}}
\end{enumerate}
\vspace{2mm}

\noindent\textbf{Q60. If f ( x ) = cos ( x ) is expanded as a half-range sine series on ( 0 , $\pi$ ) , the coefficients b n are given by b n = $\pi$ 2 $\int$ 0 $\pi$ cos ( x ) sin ( n x ) d x . What is the value of b 1 ?} \hfill [Tag: half range series | Hard]
\begin{enumerate}[label=(\alph*),topsep=2pt,itemsep=1pt,partopsep=0pt,parsep=0pt]
    \item {{1}}
    \item {{$\pi$}}
    \item {{2}}
    \item {{0}}
\end{enumerate}
\vspace{2mm}


\section{Section B: Comprehensive Subjective Problem with Detailed Analytical Rubric}

\begin{subjectiveq}{4}{1}
Obtain the Fourier series expansion of $f(x) = x^2$ in the interval $[-\pi, \pi]$. Hence deduce that $\sum_{n=1}^\infty \frac{1}{n^2} = \frac{\pi^2}{6}$.
\end{subjectiveq}

\begin{solbox}
Since $f(x) = x^2$ is an even function, $b_n = 0$. Compute $a_0 = \frac{2}{\pi}\int_0^\pi x^2 dx = \frac{2}{\pi}\left[\frac{x^3}{3}\right]_0^\pi = \frac{2\pi^2}{3}$. Compute $a_n = \frac{2}{\pi}\int_0^\pi x^2 \cos nx dx = \frac{4(-1)^n}{n^2}$ via integration by parts. Thus $x^2 = \frac{\pi^2}{3} + 4\sum_{n=1}^\infty \frac{(-1)^n}{n^2} \cos nx$. Setting $x = \pi$: $\pi^2 = \frac{\pi^2}{3} + 4\sum_{n=1}^\infty \frac{1}{n^2} \implies \sum_{n=1}^\infty \frac{1}{n^2} = \frac{\pi^2}{6}$.
\end{solbox}

\begin{examinertip}
Step 1: Identify the foundational definition and state governing theorems.\\
Step 2: Clearly write each algebraic transformation and boundary substitution.\\
Step 3: State the final numerical/algebraic invariant clearly in a concluding box.
\end{examinertip}
